\def\UseOption{hibyr1native}
\edef\UseOption{\UseOption,HIBY_R3PROII_PAD}
\edef\UseOption{\UseOption,HAVE_BACKLIGHT}
\edef\UseOption{\UseOption,HAVE_RB_BL_IN_FLASH}
\edef\UseOption{\UseOption,hibyr1native}

\newcommand{\playerman}{HiBy}
\newcommand{\playertype}{R1}
\newcommand{\playerlongtype}{\playertype}
\newcommand{\dapdisplaywidth}{480}
\newcommand{\dapdisplayheight}{800}
\newcommand{\dapdisplaydepth}{16}
\newcommand{\specimg}{hibyr1}
%Used to name the player, e.g. ...to the \dap
\newcommand{\dap}{player}
%For use when referring to the player. E.g. the \daps\ capacity ...
\newcommand{\daps}{player's}
\newcommand{\firmwarefilename}{\fname{rockbox.r1}}
\newcommand{\bootfilename}{\fname{bootloader.r1}}
\newcommand{\firmwareextension}{\fname{r1}}
\newcommand{\bootbackupfilename}{\fname{hibyr1-boot.bin}}
\newcommand{\jztoolsubcommand}{hibyr1}
\newcommand{\screenshotsize}{6cm}
\newcommand{\disk}{SD card}

% link external keymap files. The touchscreen one is not optional here: the R1
% has a CST8xx panel, apps/features.txt therefore sets the "touchscreen" option,
% and the plugin chapters reference \TouchCenter and the rest under \opt{}.
% Derived from apps/keymaps/keymap-hibyr3proii.c, which is the keymap this
% target compiles (CONFIG_KEYPAD = HIBY_R3PROII_PAD).
%
% The R1 has FIVE physical keys -- Power, Play, Volume Up, Volume Down, Next --
% and a touchscreen. It has no Prev key: button-target.h defines BUTTON_PREV
% only so the shared R3 Pro II keymap compiles, and button_read_device() can
% never report it. Actions bound solely to BUTTON_PREV are therefore reachable
% only through the touchscreen, and are documented as such below rather than
% naming a key the player does not have.

\newcommand{\ButtonPower}{\btnfnt{Power}}
\newcommand{\ButtonPlay}{\btnfnt{Play}}
\newcommand{\ButtonVolUp}{\btnfnt{Volume Up}}
\newcommand{\ButtonVolDown}{\btnfnt{Volume Down}}
\newcommand{\ButtonNext}{\btnfnt{Next}}
\newcommand{\ButtonTouch}{\btnfnt{Touchscreen}}
\newcommand{\ButtonSelect}{\ButtonPlay}
\newcommand{\ButtonBack}{\ButtonPower}
\newcommand{\ButtonMenu}{\ButtonPower}
\newcommand{\ButtonUp}{\ButtonNext}
\newcommand{\ButtonDown}{\ButtonTouch}
\newcommand{\ButtonLeft}{\ButtonVolDown}
\newcommand{\ButtonRight}{\ButtonVolUp}
\newcommand{\ButtonPrev}{\ButtonTouch}
\newcommand{\ButtonScrollFwd}{\ButtonVolUp}
\newcommand{\ButtonScrollBack}{\ButtonVolDown}

%Button actions, standard context
\newcommand{\ActionStdPrev}{\TouchActionStdPrev}
\newcommand{\ActionStdPrevRepeat}{\TouchActionStdPrevRepeat}
\newcommand{\ActionStdNext}{\ButtonNext}
\newcommand{\ActionStdNextRepeat}{Long \ButtonNext}
\newcommand{\ActionStdOk}{\ButtonPlay}
\newcommand{\ActionStdCancel}{\ButtonPower}
\newcommand{\ActionStdContext}{Long \ButtonPlay}
\newcommand{\ActionStdMenu}{\TouchActionStdMenu}
\newcommand{\ActionStdKeylock}{\ButtonPower{} + \ButtonVolUp}
\newcommand{\ActionStdQuickScreen}{\TouchActionStdQuickScreen}
\newcommand{\ActionStdHotkey}{\ButtonTouch}
\newcommand{\ActionStdUsbCharge}{\ButtonVolUp{} or \ButtonVolDown}

%Button actions, wps context
\newcommand{\ActionWpsPlay}{\ButtonPlay}
\newcommand{\ActionWpsStop}{Long \ButtonPlay}
\newcommand{\ActionWpsSkipNext}{\ButtonNext}
\newcommand{\ActionWpsSkipPrev}{\TouchActionWpsSkipPrev}
\newcommand{\ActionWpsSeekFwd}{Long \ButtonNext}
\newcommand{\ActionWpsSeekBack}{\TouchActionWpsSeekBack}
\newcommand{\ActionWpsVolUp}{\ButtonVolUp}
\newcommand{\ActionWpsVolDown}{\ButtonVolDown}
\newcommand{\ActionWpsMenu}{\ButtonPower}
\newcommand{\ActionWpsBrowse}{\TouchActionWpsBrowse}
\newcommand{\ActionWpsContext}{\TouchActionWpsContext}
\newcommand{\ActionWpsQuickScreen}{\TouchActionWpsQuickScreen}
\newcommand{\ActionWpsHotkey}{\TouchBottomRight}

%Button actions, settings context
\newcommand{\ActionSettingsInc}{\ButtonVolUp}
\newcommand{\ActionSettingsIncBigStep}{Long \ButtonVolUp}
\newcommand{\ActionSettingsDec}{\ButtonVolDown}
\newcommand{\ActionSettingsDecBigStep}{Long \ButtonVolDown}

%Button actions, list context
\newcommand{\ActionListVolUp}{\ButtonVolUp}
\newcommand{\ActionListVolDown}{\ButtonVolDown}

%Button actions, tree context
\newcommand{\ActionTreeWps}{\TouchActionTreeWps}
\newcommand{\ActionTreeEnter}{\ButtonPlay}
\newcommand{\ActionTreeParentDirectory}{\ButtonPower}
\newcommand{\ActionTreeMenu}{\ButtonTouch}
\newcommand{\ActionTreeHotkey}{\ButtonTouch}

%Button actions, yesno context
\newcommand{\ActionYesNoAccept}{\ButtonPlay}

%Button actions, quickscreen context
\newcommand{\ActionQuickScreenTop}{\ButtonVolUp}
\newcommand{\ActionQuickScreenDown}{\ButtonVolDown}
\newcommand{\ActionQuickScreenLeft}{\ButtonPlay}
\newcommand{\ActionQuickScreenRight}{\ButtonNext}
\newcommand{\ActionQuickScreenExit}{\ButtonPower}

%Button actions, pitchscreen context
\newcommand{\ActionPsIncSmall}{\ButtonVolUp}
\newcommand{\ActionPsDecSmall}{\ButtonVolDown}
\newcommand{\ActionPsIncBig}{\TouchActionPsIncBig}
\newcommand{\ActionPsDecBig}{\TouchActionPsDecBig}
\newcommand{\ActionPsNudgeLeft}{\TouchActionPsNudgeLeft}
\newcommand{\ActionPsNudgeRight}{\TouchActionPsNudgeRight}
\newcommand{\ActionPsToggleMode}{\TouchActionPsToggleMode}
\newcommand{\ActionPsReset}{\TouchActionPsReset}
\newcommand{\ActionPsExit}{\ButtonPower}
\newcommand{\ActionPsSlower}{\TouchActionPsSlower}
\newcommand{\ActionPsFaster}{\TouchActionPsFaster}

%Button actions, keyboard context
\newcommand{\ActionKbdLeft}{\ButtonVolDown}
\newcommand{\ActionKbdRight}{\ButtonVolUp}
\newcommand{\ActionKbdCursorLeft}{\TouchActionKbdCursorLeft}
\newcommand{\ActionKbdCursorRight}{\TouchActionKbdCursorRight}
\newcommand{\ActionKbdUp}{\ButtonNext}
\newcommand{\ActionKbdDown}{\TouchActionKbdDown}
\newcommand{\ActionKbdPageFlip}{\TouchActionKbdPageFlip}
\newcommand{\ActionKbdBackSpace}{\TouchActionKbdBackSpace}
\newcommand{\ActionKbdSelect}{\ButtonPlay}
\newcommand{\ActionKbdDone}{Long \ButtonPlay}
\newcommand{\ActionKbdAbort}{\ButtonPower}

%Button actions, bookmark context
\newcommand{\ActionBmDelete}{\TouchActionBmDelete}

% Stopping from the tree is bound only in keymap-touchscreen.c.
\newcommand{\ActionTreeStop}{\TouchActionTreeStop}

% Neither keymap binds A-B actions; these touchscreen aliases are placeholders.
\newcommand{\ActionWpsAbReset}{\TouchActionWpsAbReset}
\newcommand{\ActionWpsAbSetAPrevDir}{\TouchActionWpsAbSetAPrevDir}
\newcommand{\ActionWpsAbSetBNextDir}{\TouchActionWpsAbSetBNextDir}

%Plugin lib actions
\newcommand{\PluginUp}{\TouchTopMiddle}
\newcommand{\PluginDown}{\TouchBottomMiddle}
\newcommand{\PluginLeft}{\TouchMidLeft}
\newcommand{\PluginRight}{\TouchMidRight}
\newcommand{\PluginSelect}{\TouchCenter}
\newcommand{\PluginSelectRepeat}{Long \TouchCenter}
\newcommand{\PluginCancel}{\TouchBottomRight}
\newcommand{\PluginExit}{\ButtonPower}

%Bootloader actions
\newcommand{\ActionIngenicUSBBoot}{\ButtonNext}
\newcommand{\ActionBootRecoveryMenu}{\ButtonNext}
\newcommand{\ActionBootOFPlayer}{\ButtonPlay}

\input{platform/keymap-touchscreen.tex}
